\section*{अध्याय \ref{ch:b3:trace-element-modelling} --- अनुरेखी तत्त्व मॉडलिंग}

\begin{solution}{exo:b3:trace-element-modelling:1}
स्पिनेल लर्ज़ोलाइट: $0.12\times0.40 + 0.88\times0.016 = 0.048 + 0.0141 =
0.062$। गार्नेट लर्ज़ोलाइट: $0.195\times0.40 + 0.167\times0.11 +
0.639\times0.016 = 0.078 + 0.0184 + 0.0102 = 0.107$।
\end{solution}

\begin{solution}{exo:b3:trace-element-modelling:2}
$1/(0.05 + 0.02\times0.95) = 1/0.069 = \qty{14.5}{ppm}$;
$1/(0.05 + 0.10\times0.95) = 1/0.145 = \qty{6.9}{ppm}$।
\end{solution}

\begin{solution}{exo:b3:trace-element-modelling:3}
जब $F \to 0$ की ओर जाता है, तब संवर्धन सबसे अधिक होता है: $1/D = 100$। $D = 0$ होने पर,
$C_L/C_0 = 1/F = 20$।
\end{solution}

\begin{solution}{exo:b3:trace-element-modelling:4}
$\mathrm{La/Yb} = 30/2.5 = 12.0$; $(\mathrm{La}/\mathrm{Yb})_N =
(30/0.232)/(2.5/0.163) = 129.3/15.34 = 8.4$।
\end{solution}

\begin{solution}{exo:b3:trace-element-modelling:5}
बैच: $C_S/C_0 = 0.01/(0.01 + 0.05\times0.99) = 0.168$। भिन्नात्मक:
$(0.95)^{99} = 0.0062$, 27 गुना कम। तात्कालिक गलित में
$0.0062/0.01 = 0.62\,C_0$ सांद्रण रखता है: स्रोत से कम।
\end{solution}

\begin{solution}{exo:b3:trace-element-modelling:6}
$D_0 = 0.58\times0.0003 + 0.27\times0.002 + 0.12\times0.053 + 0.03\times0.01
= 0.00737$; $P = 0.20\times0.0003 + 0.25\times0.002 + 0.45\times0.053 +
0.10\times0.01 = 0.0254$। $F = 0.02$ पर अ-मोडल $1/(0.00737 +
0.02\times0.9746) = 37.2$, और मोडल $1/(0.00737 + 0.02\times0.9926) = 36.7$।
$F = 0.10$ पर: 9.54 और 9.38, जो \qty{1.7}{\percent} के अंतर पर हैं।
\end{solution}

\begin{solution}{exo:b3:trace-element-modelling:7}
$D = 0.1$: बैच $1/(0.1 + 0.09) = 5.26$; समुच्चित $(1 - 0.9^{10})/0.1 =
6.51$। $D = 0.01$: बैच $1/(0.01 + 0.099) = 9.17$; समुच्चित
$(1 - 0.9^{100})/0.1 = 10.00$। $D = 0.1$ वाला युग्म अधिक भिन्न है:
$D = 0.01$ के लिए दोनों गलित पहले ही लगभग पूरा तत्त्व रखते हैं, और दोनों
$1/F$ की ओर प्रवृत्त होते हैं।
\end{solution}

\begin{solution}{exo:b3:trace-element-modelling:8}
$D' = 0.99\times0.001 + 0.01 = 0.01099$। पहला गलित: $1/D' = 91$, \omterm{thm:b3:trace-element-modelling:fractional}{भिन्नात्मक गलन} में
1000 है। $F = 0.01$ पर: $91\times0.99^{90.0} = 36.8$, जबकि
$1000\times0.99^{999} = 0.044$: रंध्र गलित अवशेष को
तत्त्वों से वंचित होने से बचाता है।
\end{solution}

\begin{solution}{exo:b3:trace-element-modelling:9}
अनुपात $(D_{\mathrm{Sm}} + F(1 - D_{\mathrm{Sm}}))/(D_{\mathrm{La}}
+ F(1 - D_{\mathrm{La}}))$ के गुणक से बढ़ता है: $F = 0.01$ पर, $0.0714/0.0177 = 4.0$; $F = 0.10$ पर,
$0.156/0.107 = 1.46$। अनुपात छोटे $F$ पर $F$ के प्रति संवेदनशील है और
लगभग \qty{10}{\percent} से ऊपर प्रायः असंवेदी है; यह $D_{\mathrm{Sm}}$ पर भी निर्भर करता है,
अतः अवशिष्ट खनिजों पर।
\end{solution}

\begin{solution}{exo:b3:trace-element-modelling:10}
$\frac1F\int_0^F \frac1D(1 - x)^{1/D - 1}\,\mathrm dx = \frac1F\bigl[-(1 -
x)^{1/D}\bigr]_0^F = [1 - (1 - F)^{1/D}]/F$, अध्याय का सूत्र; यदि
$(1 - F)^{1/D} \to 0$ हो, तो यह $1/F$ गुना $C_0$ की ओर प्रवृत्त होता है। \qty{1}{\percent} के भीतर:
$(1 - F)^{100} < 0.01$, $F > 1 - 0.01^{0.01} = 0.045$, लगभग
\qty{4.5}{\percent}।
\end{solution}

\begin{solution}{exo:b3:trace-element-modelling:11}
ढाल $a = (10 - 6)/(20 - 10) = \qty{0.40}{ppm^{-1}}$; अंतःखंड $b = 10 -
0.40\times20 = 2.0$। $D_{\mathrm{Yb}} = 0.40\times0.44 = 0.176$;
$C_0^{\mathrm{La}} = 2.0\times0.44/0.824 = \qty{1.07}{ppm}$। $F_P = 1.07/20 =
\qty{5.3}{\percent}$, $F_Q = 1.07/10 = \qty{10.7}{\percent}$।
\end{solution}

\begin{solution}{exo:b3:trace-element-modelling:12}
$F = x_i/p_i$: क्लिनोपाइरॉक्सीन $0.12/0.45 = 0.27$, स्पिनेल $0.03/0.10 = 0.30$,
ऑर्थोपाइरॉक्सीन $0.27/0.25 = 1.08$, ओलिवीन $0.58/0.20 = 2.9$; क्लिनोपाइरॉक्सीन
पहले गलता है। $F = 0.2$ पर: $\bar D = (0.00737 - 0.0254\times0.2)/0.8 =
0.0029$। $F = 0.27$ से आगे $p_i$ बदल जाते हैं (गलने के लिए कोई क्लिनोपाइरॉक्सीन नहीं बचता),
और सूत्र, जो स्थिर $p_i$ मानता है, अब लागू नहीं होता।
\end{solution}

\begin{solution}{pb:b3:trace-element-modelling:1}
\textbf{1.} $D_{\mathrm{La}} = 0.12\times0.058 + 0.88\times0.001 = 0.00784$;
$D_{\mathrm{Yb}} = 0.12\times0.58 + 0.88\times0.16 = 0.210$।

\textbf{2.} $D_{\mathrm{La}} = 0.195\times0.058 + 0.167\times0.001 +
0.639\times0.001 = 0.0121$; $D_{\mathrm{Yb}} = 0.195\times0.58 +
0.167\times4.5 + 0.639\times0.16 = 0.967$। गोल करने पर: 0.0078, 0.21, 0.012, 0.97।

\textbf{3.} $D_{\mathrm{Yb}}$, 0.21 से 0.97 तक, क्योंकि गार्नेट
भारी दुर्लभ मृदा तत्त्वों को धारण करता है ($D = 4.5$)। इटर्बियम अभी भी बहुत थोड़ा असंगत है,
$0.97 < 1$।

\textbf{4.} $(0.648/0.232)/(0.441/0.163) = 2.79/2.71 = 1.03$: एक समतल प्रतिरूप
कॉन्ड्राइट से लगभग 2.8 गुना।

\textbf{5.} $0.648/(0.0078 + 0.01\times0.9922) = \qty{36.6}{ppm}$;
\qty{5}{\percent} पर \qty{11.3}{ppm}; \qty{10}{\percent} पर \qty{6.05}{ppm}।

\textbf{6.} $F = (0.648/9.6 - 0.0078)/0.9922 = (0.0675 - 0.0078)/0.9922 =
0.0602$, \qty{6.0}{\percent}।

\textbf{7.} $\mathrm{Yb} = 0.441/(0.21 + 0.0602\times0.79) = \qty{1.71}{ppm}$
और La/Yb $= 5.6$; लावा A में 1.7 और 5.6 हैं: यह स्पिनेल स्रोत से मेल खाता है।

\textbf{8.} La $= 0.0078\times9.6 = \qty{0.075}{ppm}$, Yb $= 0.21\times1.71 =
\qty{0.36}{ppm}$। लैंथेनम शेष: $0.940\times0.075/0.648 = 0.11$, लगभग
\qty{11}{\percent}।

\textbf{9.} $F = 0.648/9.6 = \qty{6.75}{\percent}$.

\textbf{10.} $(0.648/0.0078)(1 - F)^{127.2}$: \qty{0.1}{\percent} पर \qty{73}{ppm},
\qty{1}{\percent} पर \qty{23}{ppm},
\qty{6}{\percent} पर \qty{0.03}{ppm}।

\textbf{11.} $0.648\times(1 - 0.940^{128.2})/0.060 = \qty{10.8}{ppm}$.

\textbf{12.} $F = \qty{6.75}{\percent}$: इस $F$ पर, $(1 - F)^{1/D} =
\num{1.3e-4}$, अतः समुच्चित गलित में स्रोत का समस्त लैंथेनम होता है, और
$\bar C_L = C_0/F$, जैसा पूर्णतः असंगत तत्त्व के लिए होता है।

\textbf{13.} $0.648\times0.9325^{127.2} = \qty{9e-5}{ppm}$, यह लगभग 800 गुना
\omterm{thm:b3:trace-element-modelling:batch}{बैच गलन} के \qty{0.075}{ppm} से कम है।

\textbf{14.} दोनों मॉडल उत्तर की सीमा तय करते हैं: \qty{6.0}{\percent} (बैच)
से \qty{6.75}{\percent} (समुच्चित भिन्नात्मक)। मॉडल चाहे जो हो, लैंथेनम $F$ को
गलन के लगभग \qty{1}{\percent} के भीतर निश्चित कर देता है।

\textbf{15.} लावा A: La/Yb $= 5.6$, $(\mathrm{La}/\mathrm{Yb})_N =
(9.6/0.232)/(1.7/0.163) = 4.0$। लावा B: La/Yb $= 34.7$, $(\mathrm{La}/\mathrm{Yb})_N
= (15.6/0.232)/(0.45/0.163) = 24.4$।

\textbf{16.} जब $F \to 0$: $(0.648/0.0078)/(0.441/0.21) = 83.1/2.10 = 39.6$।

\textbf{17.} $F = \qty{0.11}{\percent}$ पर, जहाँ गलित में
\qty{72}{ppm} La और \qty{2.1}{ppm} Yb होता। लावा B में \qty{15.6}{ppm} La है:
स्पिनेल स्रोत उसे नहीं बना सकता।

\textbf{18.} $F = (0.648/15.6 - 0.012)/0.988 = 0.0299$, \qty{3.0}{\percent};
$\mathrm{Yb} = 0.441/(0.97 + 0.0299\times0.03) = \qty{0.454}{ppm}$, जैसा
देखा गया।

\textbf{19.} गार्नेट लगभग 60 से \qty{90}{km} गहराई पर स्पिनेल की जगह लेता है: लावा B
उससे अधिक गहराई पर गलना शुरू हुआ, लावा A कम गहराई पर, स्पिनेल क्षेत्र में। गार्नेट-लर्ज़ोलाइट
गलित, जिसमें \qty{9.6}{ppm} La है ($F \approx \qty{5.6}{\percent}$), उसमें La/Yb लगभग 21 होता, 5.6 नहीं:
लावा A गार्नेट-लर्ज़ोलाइट से नहीं आया।

\textbf{20.} $a = 0.97/0.441 = \qty{2.20}{ppm^{-1}}$, $b = 0.648\times0.03/0.441
= 0.044$; La $= 15.6$ पर, La/Yb $= 2.20\times15.6 + 0.044 = 34.4$, जो
लावा B के लिए 34.7 से मेल खाता है।

\textbf{21.} ओलिवीन दोनों तत्त्वों को नहीं ग्रहण करता: La/Yb अपरिवर्तित रहता है। प्राथमिक
गलित में $9.6\times0.9 = \qty{8.64}{ppm}$ La था, अतः $F = (0.648/8.64 -
0.0078)/0.9922 = \qty{6.8}{\percent}$ आता है, 6.0 के बजाय।

\textbf{22.} एक ही स्रोत दोनों लावाओं को बना सकता है: लावा B एक छोटा गलित है
(\qty{3}{\percent}) जो लगभग 60 से
\qty{90}{km} से अधिक गहराई पर, गार्नेट क्षेत्र में शुरू हुआ; लावा A स्पिनेल क्षेत्र में एक बड़ा गलित है।

\textbf{23.} लावा A स्पिनेल लर्ज़ोलाइट के \textbf{\qty[text-series-to-math=true]{6.0}{\percent}}
\omterm{thm:b3:trace-element-modelling:batch}{बैच गलन} से बना गलित है (यदि गलित को भिन्नात्मक रूप से निकाला गया
और समुच्चित किया गया हो, तो \qty{6.75}{\percent})।
\end{solution}
